\section*{APÉNDICES}
\addcontentsline{toc}{section}{APÉNDICES}

Los presentes Apéndices compilan la documentación técnica de soporte, especificaciones de la API REST, esquemas de bases de datos y fragmentos de código algorítmico fundamentales conforme a los lineamientos de reproducibilidad del estándar APA 7.

\subsection*{Apéndice A: Especificación de Endpoints y Contratos de la API REST}

El microservicio backend en Flask expone un conjunto de endpoints RESTful bajo formato JSON para consumo desde el frontend o clientes externos:

\begin{enumerate}
    \item \textbf{Inferencia Probabilística de Riesgo Espacial GMM:}
    \begin{itemize}
        \item \textbf{Ruta:} \texttt{POST /api/predict\_risk}
        \item \textbf{Cuerpo de la Petición (\textit{Request Body}):}
\begin{lstlisting}[language=json]
{
  "latitude": 41.8781,
  "longitude": -87.6298
}
\end{lstlisting}
        \item \textbf{Respuesta del Servidor (\textit{Response Body}):}
\begin{lstlisting}[language=json]
{
  "status": "success",
  "log_likelihood": -4.215,
  "risk_score_normalized": 0.842,
  "dominant_component": 1,
  "component_responsibilities": [0.72, 0.15, 0.08, 0.05],
  "risk_level": "Alto"
}
\end{lstlisting}
    \end{itemize}

    \item \textbf{Clasificación Semántica de Texto en Tiempo Real (NLP):}
    \begin{itemize}
        \item \textbf{Ruta:} \texttt{POST /api/classify\_text}
        \item \textbf{Cuerpo de la Petición (\textit{Request Body}):}
\begin{lstlisting}[language=json]
{
  "text": "Sujeto armado desciende de motocicleta y sustrae pertenencias a transeunte con arma de fuego en via publica"
}
\end{lstlisting}
        \item \textbf{Respuesta del Servidor (\textit{Response Body}):}
\begin{lstlisting}[language=json]
{
  "status": "success",
  "predicted_cluster": 3,
  "archetype_name": "Asalto a Mano Armada en Via Publica",
  "top_keywords": ["arma", "fuego", "calle", "atraco", "victima"],
  "inference_time_ms": 3.8
}
\end{lstlisting}
    \end{itemize}

    \item \textbf{Ejecución de Clustering Espacial Dinámico:}
    \begin{itemize}
        \item \textbf{Ruta:} \texttt{POST /api/run\_clustering}
        \item \textbf{Parámetros:} \texttt{radio\_km} (float), \texttt{min\_pts} (int), \texttt{algorithm} (enum: dbscan, hdbscan, optics).
    \end{itemize}
\end{enumerate}

\subsection*{Apéndice B: Implementación Algorítmica Central (Python)}

El Listado \ref{lst:haversine_dbscan} reproduce la implementación del cálculo geodésico y clustering de densidad espacial sobre coordenadas angulares en radianes:

\begin{lstlisting}[language=Python, caption={Cálculo geodésico y clustering espacial con Haversine en radianes.}, label={lst:haversine_dbscan}]
import numpy as np
import pandas as pd
from sklearn.cluster import DBSCAN, OPTICS

RADIO_TIERRA_KM = 6371.0088

def run_spatial_density(df_points, radio_km=2.5, min_pts=2):
    """
    Ejecuta DBSCAN y OPTICS utilizando la metrica geodesica de Haversine.
    Convierte previamente latitud y longitud a radianes.
    """
    p = df_points.dropna(subset=['lat', 'lon']).copy()
    if len(p) == 0:
        return p

    # Agrupacion de coordenadas identicas para eficiencia computacional
    u = p.groupby(['lat', 'lon']).size().reset_index(name='n')
    xu = np.radians(u[['lat', 'lon']].to_numpy())

    # Parametrizacion del radio en radianes
    eps_rad = radio_km / RADIO_TIERRA_KM
    db = DBSCAN(eps=eps_rad, min_samples=min_pts, metric='haversine')
    u['dbscan'] = db.fit(xu, sample_weight=u['n'].to_numpy()).labels_

    # Clustering OPTICS con ordenamiento de alcanzabilidad
    opt = OPTICS(min_samples=min_pts, metric='haversine')
    u['optics'] = opt.fit(xu).labels_

    # Reintegracion relacional al DataFrame original
    resultado = p.merge(
        u[['lat', 'lon', 'dbscan', 'optics']], 
        on=['lat', 'lon'], 
        how='left'
    )
    return resultado
\end{lstlisting}

\subsection*{Apéndice C: Arquitectura del Autoencoder Neuronal para Detección de Anomalías}

El Listado \ref{lst:autoencoder_anomalias} ilustra la red neuronal profunda de reconstrucción y determinación del umbral de anomalías:

\begin{lstlisting}[language=Python, caption={Autoencoder con arquitectura de cuello de botella y cálculo de error MSE.}, label={lst:autoencoder_anomalias}]
from sklearn.neural_network import MLPRegressor
import numpy as np

def detect_anomalies_autoencoder(X_features, percentile_threshold=90):
    """
    Entrena un autoencoder para reconstruir la matriz de caracteristicas.
    Las instancias con mayor MSE se clasifican como anomalias.
    """
    X_dense = np.asarray(X_features)

    # Autoencoder con cuello de botella D -> 16 -> 4 -> 16 -> D
    ae = MLPRegressor(
        hidden_layer_sizes=(16, 4, 16),
        activation='relu',
        max_iter=250,
        random_state=42,
        early_stopping=True
    )
    ae.fit(X_dense, X_dense)

    # Reconstruccion y error cuadratico medio por fila
    X_pred = ae.predict(X_dense)
    mse = np.mean((X_dense - X_pred) ** 2, axis=1)

    # Determinacion del umbral critico
    umbral = float(np.percentile(mse, percentile_threshold))
    es_anomalia = mse > umbral

    return mse, umbral, es_anomalia, ae
\end{lstlisting}

